% Faithful official CCC2022 article16 full Section5 main agnostic-learning proof.
\section*{5 Learning in Randomized Heuristica}
The main result of this section is the following.
\paragraph{Theorem 56.} If $\mathrm{DistNP}\subseteq\mathrm{AvgBPP}$, then for any time constructible functions $s,T,a:\mathbb N\to\mathbb N$, and $\varepsilon\in[0,1]$, $\mathrm{SIZE}[s(n)]$ is agnostic learnable on $\mathrm{Samp}[T(n)]/a(n)$ in time $\operatorname{poly}(n,\varepsilon^{-1},s(n),T(n),a(n))$ with sample complexity
\[\left(\frac{(n+s(n)+a(n)+\log T(n))^3}{\varepsilon^8}\right)^{1+o(1)}.\]
\subsection*{5.1 Ingredients}
In this section, we will use $D$ to denote a family of distributions (one for each $n$) and $\mathcal D$ to denote a class of family of distributions (e.g., those that are efficiently samplable). For a distribution $D$ and $m\in\mathbb N$, let $D^m$ denote the distribution $x_1\circ x_2\circ\cdots\circ x_m$ where $x_1,x_2,\ldots,x_m\sim D$.
\subsubsection*{5.1.1 Learning from RRHS Refutation}
\paragraph{Definition 57 (Correlative RRHS-Refutation [21, Definition 3]).} Let $\mathcal C$ be a concept class, and let $\mathcal D$ be a class of example distributions. A randomized algorithm $A$ is a correlatively random-right-hand-side-refuter (correlative RRHS-refuter) for $\mathcal C$ on $\mathcal D$ with sample complexity $m$ if $A$ satisfies the following. $A$ takes as input $n\in\mathbb N$, $\varepsilon\in(0,1)$ and a set $S=(\langle x^{(1)},b^{(1)}\rangle,\ldots,\langle x^{(m)},b^{(m)}\rangle)$ of samples, where $x^{(i)}\in\{0,1\}^n$ and $b^{(i)}\in\{0,1\}$ for every $i\in[m]$, and
\begin{itemize}
\item Soundness: if the samples $S$ are i.i.d. from a distribution $D'$ on $\{0,1\}^n\times\{0,1\}$ such that the marginal on $\{0,1\}^n$ equals $D_n$ for some $D_n\in\mathcal D_n$ and there exists $f\in\mathcal C_n$ with
\[\Pr_{\langle x^{(i)},b^{(i)}\rangle\sim D'}[b^{(i)}=f(x^{(i)})]\geq\frac12+\frac\varepsilon2,\]
then $\Pr_{S,A}[A(n,\varepsilon,S)=\text{correlative}]\geq2/3$.
\item Completeness: if the samples $S$ are selected i.i.d. so that $x^{(1)},\ldots,x^{(m)}\sim D_n$ for some $D_n\in\mathcal D_n$ and $b^{(1)},\ldots,b^{(m)}\sim U$, then
\[\Pr_{S,A}[A(n,\varepsilon,S)=\text{random}]\geq2/3.\]
\end{itemize}
\paragraph{Theorem 58 (Agnostic Learning from RRHS-Refutation [21]; see also [13, Theorem 5]).} Let $\mathcal C$ be a concept class, and let $\mathcal D$ be a class of example distributions. If there exists a correlative RRHS-refuter for $\mathcal C$ on $\mathcal D$ with sample complexity $m(n,\varepsilon)$ and running time $T(n,\varepsilon)$, then $\mathcal C$ is agnostic learnable with
\begin{align*}
\text{sample complexity}&=O\left(\frac{m(n,\varepsilon/2)^3}{\varepsilon^2}\right)\text{ and}\\
\text{running time}&=O\left(T(n,\varepsilon/2)\cdot\frac{m(n,\varepsilon/2)^2}{\varepsilon^2}\right).
\end{align*}
\subsubsection*{5.1.2 Probabilistic sampling depth}
\paragraph{Definition 59 (Probabilistic Sampling Depth).} Let $t,t'\in\mathbb N$, where $t'>t$. For a family $D=\{D_n\}_{n\geq1}$ of distributions, we introduce the $(t,t')$-probabilistic-sampling-depth functions $\operatorname{psd}_D^{t,t'}:=\{\operatorname{psd}_{D,n}^{t,t'}\}_{n\geq1}$, where
\[\operatorname{psd}_{D,n}^{t,t'}(m)=\mathbb E_{X\sim D_n^m}[\mathrm{pK}^t(X)-\mathrm{pK}^{t'}(X)].\]
For a collection $\mathcal D$ of families of distributions, we define $\operatorname{psd}_{\mathcal D}^{t,t'}:=\{\operatorname{psd}_{\mathcal D,n}^{t,t'}\}_{n\geq1}$, where
\[\operatorname{psd}_{\mathcal D,n}^{t,t'}(m)=\max_{D\in\mathcal D}\operatorname{psd}_{D,n}^{t,t'}(m).\]
\paragraph{Lemma 60 (Small Probabilistic Sampling Depth for Samplable Distributions).} There exists a polynomial $p_1:\mathbb N\times\mathbb N\to\mathbb N$ such that for any $T,a:\mathbb N\to\mathbb N$ and $n,m\in\mathbb N$, the following holds. For every $t\geq p_1(T(n),m)$, and every $t'>t$,
\[\operatorname{psd}_{\mathrm{Samp}[T(n)]/a(n),n}^{t,t'}(m)\leq O(\log m+\log T(n)+a(n)+\log t').\]
\paragraph{Proof.} The proof is essentially the same as that of [13, Lemma 6], but uses the unconditional coding theorem for probabilistic Kolmogorov complexity (Lemma 25).
Fix any $D_n\in\mathrm{Samp}[T(n)]/a(n)$. Let $p_1$ be the polynomial $p$ in Lemma 25. We have for every $t\geq p_1(T(n),m)$,
\begin{align*}
\mathbb E_{x\sim D_n^m}[\mathrm{pK}^t(x)]&\leq\mathbb E_{x\sim D_n^m}[\mathrm{pK}^{p_1(T(n),m)}(x)]\\
&\leq\mathbb E_{x\sim D_n^m}[\log(1/D_n^m(x))]+O(\log(m)+\log T(n)+a(n))\quad\text{(Lemma25)}\\
&=H(D_n^m)+O(\log(m)+\log T(n)+a(n))\\
&\leq\mathbb E_{x\sim D_n^m}[\mathrm K(x)]+O(\log(m)+\log T(n))+a(n)\\
&\leq\mathbb E_{x\sim D_n^m}[\mathrm{pK}^{t'}(x)]+O(\log(m)+\log T(n)+a(n)+\log t')\quad\text{(Lemma18)},
\end{align*}
where the second last inequality uses the fact that for any distribution $D$, its Shannon entropy $H(D)\leq\mathbb E_{y\sim D}[\mathrm K(y)]$ (see [27, Theorem 8.1.1] and [13, Lemma 5]). Rearranging the above, we get
\[\mathbb E_{x\sim D_n^m}[\mathrm{pK}^t(x)-\mathrm{pK}^{t'}(x)]\leq O(\log(m)+\log T(n)+a(n)+\log t'),\]
as desired.\hfill$\triangleleft$
\subsection*{5.2 RRHS refuters from probabilistic average-case easiness}
We first prove the following technical theorem.
\paragraph{Theorem 61.} If $\mathrm{DistNP}\subseteq\mathrm{AvgBPP}$, then for any time constructible functions $s,T,a:\mathbb N\to\mathbb N$ and for any constant $\zeta>0$, there is a correlative RRHS-refuter for $\mathrm{SIZE}[s(n)]$ under $\mathrm{Samp}[T(n)]/a(n)$ with sample complexity
\[m:=\left(\frac{n+s(n)+a(n)+\log T(n)}{\varepsilon^2}\right)^{1+\zeta}\]
and running time $\operatorname{poly}(n,m,T(n))=\operatorname{poly}(n,1/\varepsilon,T(n),a(n),s(n))$.
\paragraph{Proof.} The proof closely follows that of [13, Theorem 8].
\paragraph{The (correlative) RRHS-refuter $R$.} Let $\tau$ be the polynomial from Lemma 28. Let $\ell_s(n)\leq O(s(n)\log s(n))$ denote the number of bits needed to encode a function $f\in\mathrm{SIZE}[s(n)]$. On input $n\in\mathbb N$, $\varepsilon>0$, and a set $S=(\langle x^{(1)},b^{(1)}\rangle,\ldots,\langle x^{(m)},b^{(m)}\rangle)$ of samples, $R$ operates as follows.
\begin{enumerate}
\item Compute $t:=\max\{p_0(n\cdot m^2),p_1(T(n),m)\}$, where $p_0$ is the polynomial from Lemma 26 and $p_1$ is the polynomial from Lemma 60. Also compute $t':=t'(t)$, where $t'$ is some polynomial specified later (in Claim 62).
\item Compute
\[\beta:=\mathrm{Approx}_\tau\text{-}\mathrm{pK}(X,1^t),\text{ and }\beta':=\mathrm{Approx}_\tau\text{-}\mathrm{pK}(X\circ b,1^{t'}),\]
where $X:=x^{(1)}\circ\cdots\circ x^{(m)}$, $b:=b^{(1)}\circ\cdots\circ b^{(m)}$, and $\mathrm{Approx}_\tau\text{-}\mathrm{pK}$ is the randomized algorithm from Lemma 28.
\item Output ``correlative'' if $\beta'-\beta\leq\theta$, where $\theta:=m(1-\varepsilon^2/8)+\ell_s(n)+n+\log\tau(t)$, and output ``random'' otherwise.
\end{enumerate}
It is easy to verify that the running time of $R$ is $\operatorname{poly}(n,m,T(n))$. Next, we argue its correctness.
\paragraph{Soundness.} Suppose we are in the ``correlative'' case. That is, the labelled examples in $S$ are sampled i.i.d from some distribution $D'$ on $\{0,1\}^n\times\{0,1\}$, whose marginal on $\{0,1\}^n$ is given by some $D\in\mathrm{Samp}[T(n)]/a(n)$, and there exists $f\in\mathrm{SIZE}[s(n)]$ such that
\[\Pr_{\langle x^{(i)},b^{(i)}\rangle\sim D'}[b^{(i)}=f(x^{(i)})]\geq\frac12+\frac\varepsilon2.\]
In this case, by a standard concentration bound, the probability over $S\sim(D')^m$ that
\[|\{i\in[m]\mid b_i=f(x^{(i)})\}|<(1/2+\varepsilon/4)\cdot m\]
is at most $\exp(-2m(\varepsilon/4)^2)\leq o(1)$, where the last inequality relies on our choice of $m$. Now observe the following.
\paragraph{Claim 62.} There exists a polynomial $t'$ such that for any $b\in\{0,1\}^m$ satisfying
\[|\{i\in[m]\mid b_i=f(x^{(i)})\}|\geq(1/2+\varepsilon/4)\cdot m,\]
we have
\[\mathrm{pK}^{t'(t)}(X\circ b)\leq\mathrm{pK}^{\tau(t)}(X)+\ell_s(n)+(1-\varepsilon^2/8)\cdot m.\]
\paragraph{Proof of Claim 62.} Note that given $X$, we can compute $f(x^{(1)}),\ldots,f(x^{(m)})$ in time $\operatorname{poly}(m\cdot\ell_s(n))$ using the encoding of $f$, which is of $\ell_s(n)$ bits. Note that $b$ and $f(x^{(1)}),\ldots,f(x^{(m)})$ disagree on at most $(1/2-\varepsilon/4)\cdot m$ coordinates. Then to recover $b$, we can define a string $e\in\{0,1\}^m$ such that $e_i=1$ iff $f(x^{(i)})\neq b_i$. Note that $e$ has hamming weight at most $(1/2-\varepsilon/4)\cdot m$. Using the inequality
\[\sum_{i=0}^k\binom mi\leq2^{H_2(k/m)\cdot m},\]
where $H_2$ is the binary entropy function, $e$ can be encoded with $H_2(1/2+\varepsilon/4)\cdot m$ bits, by lexicographic indexing among binary strings of bounded hamming weight. Using the Taylor series of $H_2$ in the neighborhood of $1/2$, for $\delta:=\varepsilon/4$, we have
\[H_2(1/2+\delta)=1-\frac1{2\ln2}\sum_{i=1}^\infty\frac{(2\delta)^{2i}}{i(2i-1)}\leq1-\frac2{\ln2}\delta^2\leq1-2\delta^2.\]
Moreover, using such an encoding, it is not hard to see that $e$ can be reconstructed in time $\operatorname{poly}(n,m)$. Therefore, given $X$, we can compute $b$ in time $\operatorname{poly}(n,m)$ using strings of length $\ell_s(n)$ (which encodes $f$) and $(1-\varepsilon^2/8)\cdot m$ (which encodes $e$). In other words, if $X$ has probabilistic $\tau(t)$-time-bounded Kolmogorov complexity $k$, then
\[\mathrm{pK}^{t'(t)}(X\circ b)\leq k+\ell_s(n)+(1-\varepsilon^2/8)\cdot m,\]
for some $t'(t):=\operatorname{poly}(\tau(t))$. This completes the proof of the claim.\hfill$\triangleleft$
Now, suppose in Step 2 of $R$, $\beta$ and $\beta'$ output by the algorithm $\mathrm{Approx}_\tau\text{-}\mathrm{pK}$ are good approximations as given in Lemma 28. Note that this happen with high probability. Then by a union bound, with probability at least $2/3$, over the samples $S\sim(D')^m$ and the internal randomness of $R$, we have
\begin{align*}
\beta'&\leq\mathrm{pK}^{t'}(X\circ b)&&\beta'\text{ is a good approximation}\\
&\leq\mathrm{pK}^{\tau(t)}(X)+\ell_s(n)+(1-\varepsilon^2/8)\cdot m&&\text{Claim62}\\
&\leq\beta+\log\tau(t)+\ell_s(n)+(1-\varepsilon^2/8)\cdot m&&\beta\text{ is a good approximation},
\end{align*}
which implies
\[\beta'-\beta\leq(1-\varepsilon^2/8)\cdot m+\ell_s(n)+n+\log\tau(t)=\theta,\]
and $R$ will output ``correlative''.
\paragraph{Completeness.} Suppose we are in the ``random'' case. That is, $b$ is sampled from $U_m$. Assuming $\mathrm{DistNP}\subseteq\mathrm{AvgBPP}$, by combining symmetry of information (Lemma 26) and probabilistic incompressibility (Lemma 20), for any $X\in\{0,1\}^{nm}$ and some polynomial $p_w$, we have
\[\Pr_{b\sim U_m}[\mathrm{pK}^{\tau(t')}(X\circ b)\geq\mathrm{pK}^{p_w(\tau(t'))}(X)+(m-10)-\log p_w(\tau(t'))]\geq1-1/8.\tag{40}\]
Let us define $p_\tau(t):=p_w(\tau(t'(t)))$. By Markov, we have
\begin{align*}
&\Pr_{X\sim D^m}[\mathrm{pK}^t(X)-\mathrm{pK}^{p_\tau(t)}(X)>8\cdot\operatorname{psd}_{D,n}^{t,p_\tau(t)}(m)]\\
&\leq\frac{\mathbb E_{X\sim D^m}[\mathrm{pK}^t(X)-\mathrm{pK}^{p_\tau(t)}(X)]}{8\cdot\operatorname{psd}_{D,n}^{t,p_\tau(t)}(m)}=1/8.\tag{41}
\end{align*}
Again, let us assume that in Step 2 of $R$, $\beta$ and $\beta'$ output by the algorithm $\mathrm{Approx}_\tau\text{-}\mathrm{pK}$ are good approximations. Then with with probability at least $1-(1/8+1/8+o(1))\geq2/3$ over $b\sim U_m$, $X\sim D^m$ and the internal randomness of $R$, we have
\begin{align*}
\beta'&\geq\mathrm{pK}^{\tau(t')}(X\circ b)-\log\tau(t')&&\beta'\text{ is a good approximation}\\
&\geq\mathrm{pK}^{p_w(\tau(t'))}(X)+m-\log(p_w(\tau(t'))\cdot\tau(t'))&&\text{Equation40}\\
&=\mathrm{pK}^t(X)-(\mathrm{pK}^t(X)-\mathrm{pK}^{p_\tau(t)}(X))+m-10-\log(p_\tau(t)\cdot\tau(t'))\\
&\geq\beta-8\cdot\operatorname{psd}_{D,n}^{t,p_\tau(t)}(m)+m-10-\log(p_\tau(t)\cdot\tau(t'))&&\text{Equation41},
\end{align*}
which implies
\[\beta'-\beta\geq m-8\cdot\operatorname{psd}_{D,n}^{t,p_\tau(t)}(m)-\log(p_\tau(t)\cdot\tau(t'))-10.\]
Now we want the above to be greater than $\theta$. We have
\begin{align*}
(\beta'-\beta)-\theta
&\geq(m-8\cdot\operatorname{psd}_{D,n}^{t,p_\tau(t)}(m)-\log(p_\tau(t)\cdot\tau(t'))-10)\\
&\quad-(m+\ell_s(n)+n+\log\tau(t)-m\varepsilon^2/8)\\
&=m\varepsilon^2/8-(n+\ell_s(n)+8\cdot\operatorname{psd}_{D,n}^{t,p_\tau(t)}(m)\\
&\qquad+\log(p_\tau(t)\cdot\tau(t')\cdot\tau(t))+10)\\
&\geq m\varepsilon^2/8-[n+\ell_s(n)+c\cdot(\log m+\log T(n)+a(n))\\
&\qquad+\log p_2(n,m,T(n))+10],\qquad\text{(Lemma60)}
\end{align*}
where $c>0$ is some constant and $p_2$ is some polynomial (that depends on $\tau,p_0,p_1$ and $p_w$). From here, it can be easily seen that the above is greater than $0$ by our choice of $m$, provided that $n$ is sufficiently large. Hence $R$ will output ``random''. This completes the proof of the theorem.\hfill$\triangleleft$
\subsection*{5.3 Concluding the proof}
We are now ready to prove Theorem 56.
\paragraph{Theorem 63 (Reminder of Theorem 56).} If $\mathrm{DistNP}\subseteq\mathrm{AvgBPP}$, then for any time constructible functions $s,T,a:\mathbb N\to\mathbb N$, and $\varepsilon\in[0,1]$, $\mathrm{SIZE}[s(n)]$ is agnostic learnable on $\mathrm{Samp}[T(n)]/a(n)$ in time $\operatorname{poly}(n,\varepsilon^{-1},s(n),T(n),a(n))$ with sample complexity
\[\left(\frac{(n+s(n)+a(n)+\log T(n))^3}{\varepsilon^8}\right)^{1+\zeta},\]
where $\zeta>0$ is any arbitrary small constant.
\paragraph{Proof.} The theorem follows by combining Theorem 61 and Theorem 58.\hfill$\triangleleft$
